% SBEG 361 Introduction to Robotics -- Lab 4, Formal report 1 (IEEE two-column, <= 6 pages + appendices)
% Compile with:  pdflatex report_template.tex   (needs the IEEEtran class, included in TeX Live / MiKTeX;
% on Overleaf choose the "IEEE Conference Template").  See lab_handout.md for what to write in each
% section and for the 100-point rubric.
\documentclass[conference]{IEEEtran}
\usepackage{amsmath,amssymb}
\usepackage{graphicx}
\usepackage{booktabs}
\usepackage{siunitx}
\usepackage{cite}

\newcommand{\Ad}[1]{[\mathrm{Ad}_{#1}]}

\begin{document}

\title{Gravity-Compensated Force Sensing at the Tip of a Laparoscopic Instrument and at an EOD Gripper}

\author{\IEEEauthorblockN{<Name>, <Student ID>}
\IEEEauthorblockA{SBEG 361 Introduction to Robotics, Systems and Biomedical Engineering\\
Repository: <URL>, tag \texttt{lab4}}}

\maketitle

\begin{abstract}
% <= 150 words: problem, method, key numbers (identified masses and centres of mass, tip-force RMS error,
% port-force detection, tipping moment), one-sentence conclusion.
\end{abstract}

\begin{IEEEkeywords}
wrench, adjoint map, force--torque sensor, gravity compensation, <application>
\end{IEEEkeywords}

\section{Introduction}
% Why tip forces matter in laparoscopy; why the package's weight and centre of mass matter for an EOD robot.
% Objectives (3 bullets).

\section{Theory}
% Exponential coordinates on SE(3) (MR Prop. 3.25, Eq. 3.88; log: Eqs. 3.91--3.92).
% Wrench F = (m, f) (MR Eq. 3.93); power invariance (Eq. 3.94); F_b = [Ad_{T_ab}]^T F_a (Eq. 3.95).
\begin{equation}
\mathcal F_b = \Ad{T_{ab}}^{\mathsf T}\mathcal F_a
\end{equation}

\section{Method}
% Sensor convention (MR Example 3.28); linear payload model and least squares; compensation; tip transfer;
% m_t = 0 consistency test and threshold; composite centre of mass; frames figure.

\section{Verification}
% Tests against modern_robotics and MR Examples 3.26 and 3.28; one hand-calculated check.

\section{Results}
\begin{table}[t]
\caption{Identified parameters (true values from the data description)}
\label{tab:params}
\centering
\begin{tabular}{@{}lccc@{}}
\toprule
Body & Mass (kg) & Centre of mass in $\{f\}$ (m) & Error \\
\midrule
Laparoscope & & & \\
EOD gripper & & & \\
Package & & & \\
\bottomrule
\end{tabular}
\end{table}

\begin{figure}[t]
\centering
% \includegraphics[width=\columnwidth]{fig1_tip_force.pdf}
\caption{Tip force components and $\|f_t\|$ against time, with and without gravity compensation. <What to see.>}
\label{fig:tipforce}
\end{figure}
% Fig. 2: |m_t| with threshold; Fig. 3: EOD frames and line of action; Fig. 4: your choice.

\section{Discussion}
% Error sources; tip versus port ambiguity; safety (patient, operator); design improvement.

\section{Conclusion}

\begin{thebibliography}{9}
\bibitem{mr} K. M. Lynch and F. C. Park, \emph{Modern Robotics: Mechanics, Planning, and Control}.
Cambridge, U.K.: Cambridge Univ. Press, 2017 (updated preprint 2019), Sec. 3.3.3 and 3.4, pp. 102--109.
\bibitem{mls} R. M. Murray, Z. Li, and S. S. Sastry, \emph{A Mathematical Introduction to Robotic Manipulation}.
Boca Raton, FL, USA: CRC Press, 1994, Ch. 2, Sec. 5, pp. 61--66.
\bibitem{mrcode} K. M. Lynch \emph{et al.}, ``Modern Robotics software (Python),'' modern\_robotics package.
\end{thebibliography}

\appendices
\section{AI-use note}
% Tool and version, what you asked, what you kept or changed, how you verified it (or "no AI tools used").
\section{Reflection}
% <= 200 words.
\section{Reproduction}
% Commands that regenerate every number and figure, expected run time.

\end{document}
